% !TeX program = xelatex
% !TeX encoding = UTF-8
% !TeX root = main.tex

\documentclass[
  oneside,
  degree=bachelor,
  field=science,
  fullwidthstop=circle,
  fontset=windows,
  times=false,
  minted=false,
  biblatex=false,
  algo=algpseudocode,
]{tongjithesis}

\renewcommand{\tjcovertext}{课程大作业报告}
\renewcommand{\tjheadertext}{课程大作业报告}
\geometry{a4paper,top=4.0cm,bottom=2.7cm,left=2.5cm,right=2.5cm}
\tjbindinglinefalse

\newcommand{\code}[1]{\texttt{#1}}
\newcommand{\todoitem}[1]{\textbf{待补充：}#1}
\newcommand{\placeholderfigure}[2]{
  \begin{figure}[htbp]
    \centering
    \fbox{
      \parbox{0.86\linewidth}{
        \centering
        \vspace{1.1cm}
        \textbf{#1}\\[0.6em]
        #2
        \vspace{1.1cm}
      }
    }
    \caption{#1}
  \end{figure}
}

\begin{document}

\school{电子信息学院}
\major{人工智能}
\student{2352396 / 2352620 / 2351283}{禹尧珅 / 蒋文涛 / 吴凯}
\thesistitle{基于 OrangePi 的两自由度云台目标跟踪系统}{综合设计实践B 课程大作业报告}
\thesistitleeng{Two-axis Pan-tilt Target Tracking System Based on OrangePi}{Course Project Report for Comprehensive Design Practice B}
\thesisadvisor{龚炜，王晓年}
\thesisdate{2026}{5}{20}

\MakeCover
\cleardoublepage

\frontmatter

\MakeAbstract{
本报告记录了综合设计实践B课程大作业“基于 OrangePi 的两自由度云台目标跟踪系统”的完整实现方案与当前阶段成果。系统以 OrangePi 5 Pro/RK3588S 为主控平台，连接 USB 摄像头、PCA9685 PWM 舵机驱动板和两自由度云台，构建了从视频采集、目标检测、状态判断到舵机控制的闭环系统。主线采用 Python、OpenCV、HSV 颜色阈值分割、形态学去噪、轮廓筛选和双舵机限幅控制，使高对比颜色目标进入画面后能够被检测、框选并尽量保持在画面中心附近。当前版本已经完成环境检查、I2C 与 PCA9685 识别、舵机测试、摄像头读取、网页远程控制台、日志记录和一键脚本化运行流程。项目还围绕课程加分项落实了 C/C++ wheel 视觉算子加速、pymp 并行加速对比、网络控制台与 mock 云端上传、云台安全保护等功能。报告中对尚需最终实测补全的数据和图片均设置了占位，便于后续在 OrangePi 板端完成最终演示后直接替换。
}{OrangePi；两自由度云台；目标跟踪；OpenCV；PCA9685；网页控制台；C++ 加速}

\MakeAbstractEng{
This report presents a two-axis pan-tilt target tracking system developed for Comprehensive Design Practice B. The system uses an OrangePi 5 Pro/RK3588S as the main controller, with a USB camera, a PCA9685 PWM servo driver and a two-degree-of-freedom pan-tilt platform. A complete closed loop from video capture, target detection and state management to servo control has been implemented. The main pipeline is based on Python, OpenCV, HSV color thresholding, morphology-based denoising, contour filtering and bounded gimbal control. The current version supports environment checks, I2C/PCA9685 validation, servo tests, camera validation, a web-based remote control console, logging and script-based execution. Bonus features include a C/C++ wheel acceleration experiment, pymp parallel benchmark, network console, mock cloud upload and gimbal safety protection.
}{OrangePi, Pan-tilt Control, Target Tracking, OpenCV, PCA9685, Web Console}

\clearpage
\tableofcontents
\cleardoublepage

\mainmatter

\chapter{项目概述}

\section{课题背景}

摄像头云台目标跟踪系统同时涉及嵌入式硬件、图像处理、控制算法、网络调试和实验数据记录，是一个典型的软硬件结合课程设计题目。系统不仅要“看见”目标，还要让机械结构根据视觉偏差实时调整姿态，因此比单纯的图像识别程序更强调工程闭环、运行稳定性和现场可演示性。

本项目使用课程提供的 OrangePi、两自由度云台、USB 摄像头、PCA9685 舵机驱动板、舵机与连接线等硬件，搭建一个可现场演示的目标跟踪原型系统。考虑到课程验收更重视稳定运行、可解释流程和可量化结果，项目不把复杂 DNN/RKNN 检测作为当前主线，而是采用高对比颜色标记物作为基准目标，优先保证摄像头、视觉、控制、舵机和远程监控可靠跑通。

\section{项目目标}

当前版本的目标是形成一条完整、可复现、可答辩说明的演示链路：

\begin{enumerate}
  \item OrangePi 上电后可进入固定项目目录，通过统一 shell 脚本完成环境检查、硬件检查和跟踪运行。
  \item USB 摄像头接入 OrangePi，程序能够稳定读取彩色图像帧。
  \item PCA9685 通过 I2C 被识别，实际地址为 \code{0x40}，All-Call 地址为 \code{0x70}。
  \item 当前舵机通道约定为 \code{channel 0 = tilt}，\code{channel 1 = pan}。
  \item 高对比颜色目标进入画面后，系统能够显示目标框、目标中心点、画面中心线、状态、FPS、pan/tilt 角度和误差信息。
  \item 目标移动时，云台连续转动并尽量将目标拉回画面中心附近。
  \item 目标短暂丢失时，系统进入保持或小范围搜索状态，而不是无约束乱转。
  \item 运行过程自动生成日志，后续可统计平均帧率、中心误差、目标丢失事件和重捕时间。
  \item 加分项与主线系统相关联，包括 C++ wheel 加速、pymp 并行对比、网络控制台、mock 云端上传和安全保护机制。
\end{enumerate}

\section{小组成员与分工}

\begin{table}[htbp]
  \centering
  \caption{小组成员信息与当前分工}
  \label{tab:members}
  \begin{tabular}{lll}
    \toprule
    学号 & 姓名 & 当前主要工作 \\
    \midrule
    2352396 & 禹尧珅 & 项目统筹、硬件调试、代码集成、报告整理 \\
    2352620 & 蒋文涛 & 视觉检测、跟踪流程、实验测试 \\
    2351283 & 吴凯 & 舵机控制、性能测试、展示材料整理 \\
    \bottomrule
  \end{tabular}
\end{table}

\section{当前完成情况}

项目目前已经从“能否启动”推进到“可进行网页远程跟踪演示”的阶段。OrangePi 端已建立虚拟环境 \code{/home/orangepi/.venvs/zb}，项目目录统一为 \code{/home/orangepi/zb/tracker\_project}。通过 \code{scripts/run} 目录下的一键脚本，可以完成环境检查、硬件链路检查、网页模拟跟踪、网页真实跟踪、C++ benchmark、pymp benchmark、日志分析、mock 云端上传和网络检查。

当前主线建议按“网页控制台模式”进行最终演示：先运行硬件检查，再运行网页模拟模式确认摄像头和识别稳定，最后运行网页真实模式控制云台。这样可以避开 MobaXterm X11 转发 OpenCV 窗口不稳定的问题，把电脑浏览器作为显示端，把 OrangePi 作为采集、计算和控制端。

\chapter{系统总体设计}

\section{总体闭环流程}

系统总体闭环流程如图~\ref{fig:system-flow} 所示。摄像头采集图像后，视觉模块输出目标是否存在、目标中心、目标框、面积和置信度；状态机根据连续检测结果判断当前处于等待、跟踪、短时丢失或搜索状态；控制模块将目标中心相对画面中心的偏差转换为云台角度；硬件模块通过 PCA9685 输出 PWM 信号控制两个舵机；日志模块记录每帧运行数据；网页模块将叠加后的实时画面通过 MJPEG 推送给浏览器。

\begin{figure}[htbp]
  \centering
  \fbox{
    \parbox{0.9\linewidth}{
      \centering
      摄像头采集 $\rightarrow$ HSV 目标检测 $\rightarrow$ 状态机判断 $\rightarrow$ 限幅控制 $\rightarrow$ PCA9685 PWM $\rightarrow$ 双舵机云台动作\\[0.6em]
      \small 日志记录与网页控制台并行输出，用于调试、展示和后续数据分析。
    }
  }
  \caption{系统总体闭环流程}
  \label{fig:system-flow}
\end{figure}

\section{硬件连接设计}

硬件连接的核心原则是：OrangePi 负责 I2C 通信和视觉计算，PCA9685 负责输出稳定 PWM，舵机电源尽量独立供电，OrangePi 与舵机电源必须共地。当前已经验证 PCA9685 位于 \code{/dev/i2c-1}，地址为 \code{0x40}。

\begin{table}[htbp]
  \centering
  \caption{当前硬件接线表}
  \label{tab:wiring}
  \begin{tabular}{lll}
    \toprule
    模块端口 & 连接对象 & 说明 \\
    \midrule
    PCA9685 VCC & OrangePi 40Pin 3.3V & I2C 逻辑供电 \\
    PCA9685 GND & OrangePi 40Pin GND & 必须与舵机电源共地 \\
    PCA9685 SDA & OrangePi 40Pin 第 3 脚 & I2C 数据线 \\
    PCA9685 SCL & OrangePi 40Pin 第 5 脚 & I2C 时钟线 \\
    PCA9685 V+ / 绿色端子 & 独立 5V 舵机电源 & 给舵机供电，不建议由 OrangePi USB 供电 \\
    PCA9685 channel 0 & 俯仰舵机 & 控制摄像头上下方向 \\
    PCA9685 channel 1 & 水平舵机 & 控制底座左右旋转 \\
    USB 摄像头 & OrangePi USB 口 & 图像采集必须接在 OrangePi 上 \\
    OrangePi Type-C & OrangePi 电源适配器 & 给开发板供电 \\
    USB 转网口 & 电脑与 OrangePi 之间 & 后续用于更稳定的远程调试与网络共享 \\
    \bottomrule
  \end{tabular}
\end{table}

\placeholderfigure{最终实物接线图占位}{后续补充一张清晰俯拍图，标出 OrangePi、PCA9685、舵机电源、SDA、SCL、GND、channel 0 和 channel 1。}

\section{软件模块划分}

当前代码按职责拆分为多个模块，便于调试和答辩说明。

\begin{table}[htbp]
  \centering
  \caption{软件模块划分}
  \label{tab:modules}
  \footnotesize
  \setlength{\tabcolsep}{4pt}
  \begin{tabular}{p{0.14\linewidth}p{0.31\linewidth}p{0.45\linewidth}}
    \toprule
    模块 & 主要文件 & 功能 \\
    \midrule
    主入口 & \code{main.py} & 解析运行模式、配置文件、跟踪后端和网页参数 \\
    配置 & \code{config.py}，\code{configs/default\_config.json} & 管理摄像头、视觉、控制、状态机、日志和硬件参数 \\
    摄像头 & \code{camera.py} & 使用 OpenCV 读取 USB 摄像头 \\
    视觉检测 & \code{tracker.py} & HSV 阈值、形态学处理、轮廓筛选、中心自动标定和动态 HSV 自适应 \\
    控制 & \code{control.py} & 根据中心偏差输出 pan/tilt 目标角度 \\
    硬件 & \code{hardware.py} & PCA9685 寄存器控制、舵机角度映射和 mock 模式 \\
    状态机 & \code{state\_machine.py} & 管理 IDLE、TRACKING、LOST\_SHORT、SEARCH 状态 \\
    网页控制台 & \code{web\_server.py} & HTTP/MJPEG 推流、实时状态 API、中心标定、开始跟踪和安全控制 \\
    日志 & \code{logging\_utils.py} & 输出每帧 CSV 与 summary JSON，并修复 sudo 生成文件权限 \\
    脚本 & \code{scripts/run/*.sh} & 将调试、演示和加分实验流程脚本化 \\
    \bottomrule
  \end{tabular}
\end{table}

\section{运行脚本体系}

为减少手输命令出错，当前版本将关键步骤整理为 \code{scripts/run} 下的 shell 脚本。总原则是：先主线、后创新点、最后回到主线采集最终数据。

\begin{table}[htbp]
  \centering
  \caption{当前推荐运行顺序}
  \label{tab:scripts}
  \footnotesize
  \setlength{\tabcolsep}{3pt}
  \begin{tabular}{p{0.07\linewidth}p{0.35\linewidth}p{0.45\linewidth}}
    \toprule
    步骤 & 脚本 & 用途 \\
    \midrule
    00 & \code{00\_import\_project.sh} & 解压上传的新版本项目，可自动识别 \code{v*.zip} \\
    01 & \code{01\_env\_check.sh} & 检查虚拟环境、依赖、包导入和单元测试 \\
    02 & \code{02\_hardware\_check.sh} & 检查 I2C、PCA9685、舵机和摄像头 \\
    03 & \code{03\_web\_console\_simulate.sh} & 网页控制台模拟模式，不控制真实舵机 \\
    04 & \code{04\_web\_console\_real.sh} & 网页控制台真实模式，控制真实舵机 \\
    05 & \code{05\_cpp\_benchmark.sh} & 创新点1：C++ wheel 视觉算子加速对比 \\
    06 & \code{06\_pymp\_benchmark.sh} & 创新点3：pymp 并行加速对比 \\
    07 & \code{07\_mock\_cloud\_upload.sh} & 创新点4：mock 云端上传摘要和 benchmark \\
    08 & \code{08\_analyze\_latest\_logs.sh} & 分析最新跟踪日志，生成报告/PPT 可用数据 \\
    09 & \code{09\_check\_materials.sh} & 交作业前检查关键材料目录 \\
    可选 & 网络共享说明 & USB 转网口到了以后临时检查网络，不属于主线必跑 \\
    \bottomrule
  \end{tabular}
\end{table}

最终远程演示建议使用如下顺序：

\begin{verbatim}
cd ~/zb/tracker_project
bash scripts/run/01_env_check.sh
bash scripts/run/02_hardware_check.sh
bash scripts/run/03_web_console_simulate.sh
bash scripts/run/04_web_console_real.sh
\end{verbatim}

创新点验证在主线跑通后补充：

\begin{verbatim}
bash scripts/run/05_cpp_benchmark.sh
bash scripts/run/06_pymp_benchmark.sh
bash scripts/run/07_mock_cloud_upload.sh
bash scripts/run/08_analyze_latest_logs.sh
\end{verbatim}

\chapter{主线功能实现}

\section{环境与依赖}

OrangePi 端统一使用虚拟环境 \code{/home/orangepi/.venvs/zb}，避免系统 Python 与项目依赖混在一起。项目依赖包括 \code{numpy}、\code{opencv-python}、\code{smbus2}、\code{pybind11}、\code{pymp-pypi}、\code{matplotlib} 和 \code{wheel}。环境检查脚本会安装依赖、验证 \code{cv2} 与 \code{smbus2} 能导入、验证 \code{orangepi\_tracker} 源码包可导入，并运行单元测试。

\begin{figure}[htbp]
  \centering
  \includegraphics[width=0.92\linewidth]{figures/env_check_current.png}
  \caption{当前版本环境检查脚本运行结果}
  \label{fig:env-check-current}
\end{figure}

\section{硬件检查}

硬件检查脚本会列出 \code{/dev/i2c-*}，扫描 \code{/dev/i2c-1} 上的 PCA9685，运行底层舵机测试、主程序舵机测试和摄像头读取测试。当前底层 I2C 扫描脚本已修改为只有明确找到 \code{0x40} 才返回成功，避免因为其它 I2C 设备存在而误判。底层舵机测试也改为安全角度的一轮测试，不再无限循环或打到极限角。

\begin{figure}[htbp]
  \centering
  \includegraphics[width=0.92\linewidth]{figures/hardware_check_current.jpg}
  \caption{当前版本硬件检查与摄像头读取结果}
  \label{fig:hardware-check-current}
\end{figure}

\section{视觉检测}

视觉检测主流程采用传统视觉路线。程序先对 BGR 图像做轻微模糊，再转换到 HSV 空间；当前网页版本不再要求手动选择颜色，而是把目标便利贴放在画面中心后进行自动标定。标定得到的 custom HSV 范围会作为后续跟踪阈值，跟踪过程中再进行小幅、慢速的动态 HSV 自适应。阈值分割后，程序通过形态学开闭运算去除噪声点，最后根据轮廓面积、最大面积比例、旋转矩形长宽比和填充率筛选目标。

检测结果包括：是否找到目标、中心点坐标、目标框、面积、置信度和 mask 预览图。叠加显示时，绿色框表示候选框，黄色点表示目标中心，白色十字线表示画面中心。这个设计便于现场演示，也便于观察中心自动标定后的 HSV 范围是否稳定。

\section{云台控制与状态机}

控制模块以画面中心为期望位置，误差定义为：

\begin{equation}
  e_x = c_x - \frac{W}{2}, \qquad e_y = c_y - \frac{H}{2}
\end{equation}

其中 $c_x,c_y$ 为目标中心坐标，$W,H$ 为图像宽高。当前配置中，水平角范围为 $20^\circ$ 到 $160^\circ$，俯仰角范围为 $35^\circ$ 到 $145^\circ$，中心角均为 $90^\circ$。控制器设置了死区、最大单帧步进和平滑系数，用于降低抖动和防止舵机突然大幅动作。

状态机包括四种状态：\code{IDLE}、\code{TRACKING}、\code{LOST\_SHORT} 和 \code{SEARCH}。连续检测到目标后进入跟踪；目标短时丢失后保持姿态；丢失超过阈值后进入小范围搜索。当前版本已经将搜索范围限制在 \code{search\_pan\_span} 附近，不再全范围乱扫；当搜索过程中重新看到目标时，会先停止搜索，等待连续锁定后再进入跟踪，提升重捕稳定性。

\section{云台安全保护}

安全保护是当前版本的重要加分点之一，主要包括：

\begin{enumerate}
  \item \textbf{角度限位：}pan 和 tilt 均设置机械安全范围，超过范围会被夹紧。
  \item \textbf{最大步进限制：}单帧最大角度变化受限，防止舵机突然大幅动作。
  \item \textbf{误差死区：}目标接近中心时不再频繁微调，减少抖动。
  \item \textbf{中心稳定保持：}目标连续多帧接近中心且移动很小时，暂时保持当前姿态。
  \item \textbf{网页急停：}网页控制台提供 STOP 按钮，按下后舵机不再继续追踪。
  \item \textbf{开始、急停与复位：}网页控制台只保留开始跟踪、急停和渐进复位。复位会先进入急停，短暂释放 PWM，再用缓入缓出的曲线小步移动到初始角，最后再次释放 PWM。
  \item \textbf{真实硬件失败显式报错：}真实模式下若 I2C/PCA9685 初始化失败，不再悄悄切到 mock，避免误判。
  \item \textbf{日志权限修复：}sudo 运行产生的日志会在关闭时尝试归还普通用户，减少后续删除和分析时的权限问题。
\end{enumerate}

\section{网页控制台}

网页控制台是当前主线演示入口，不再只是附属监看功能。它基于 HTTP/MJPEG 实现，OrangePi 负责读取摄像头和运行跟踪算法，电脑浏览器负责显示实时画面并发送控制请求。

网页控制台包含以下功能：

\begin{enumerate}
  \item 实时视频流与目标框、中心线、状态文字叠加。
  \item 实时显示 FPS、检测耗时、控制耗时、误差、pan/tilt、目标面积和置信度。
  \item 中心自动标定，将画面中心附近 ROI 的颜色转换为 custom HSV 范围。
  \item 当前 HSV 只读显示，现场不再提供手动调参和保存按钮，避免越调越乱。
  \item 开始跟踪、急停和复位三项云台安全控制。
\end{enumerate}

早期 MJPEG 推流如果直接在浏览器请求循环中读取摄像头和处理图像，网络波动或浏览器刷新可能拖慢主循环。当前版本已改为后台 producer 线程持续生成最新 JPEG 帧，并用条件变量缓存最新帧；浏览器端只读取已经缓存好的 JPEG，因此慢网络不会直接阻塞摄像头和控制逻辑。

\begin{figure}[htbp]
  \centering
  \includegraphics[width=0.92\linewidth]{figures/web_monitor_current.png}
  \caption{当前版本网页控制台效果}
  \label{fig:web-monitor-current}
\end{figure}

\section{网页控制台展示材料占位}

网页控制台后续会成为最终答辩时最直观的展示入口，因此报告中需要预留几类截图位置。最终补图时建议统一使用电脑浏览器截图，画面中保留 OrangePi 的访问地址、实时视频流、状态面板和控制按钮，这样能够同时说明“远程访问”“实时跟踪”和“可交互控制”三个信息点。

\placeholderfigure{中心自动标定前截图占位}{后续补充目标便利贴放在画面中心、尚未开始跟踪时的网页截图。}

\placeholderfigure{自动标定后 HSV 只读显示截图占位}{后续补充点击自动标定后的网页截图，重点展示 custom HSV 参数、mask 预览和目标框。}

\placeholderfigure{中心自动标定截图占位}{后续补充点击自动标定按钮后的截图，画面中心区域对准目标颜色，并展示 custom HSV 参数已经被更新。}

\placeholderfigure{网页实时状态面板截图占位}{后续补充 FPS、state、err\_x、err\_y、pan、tilt、area 和 confidence 等状态字段清晰可见的截图。}

\chapter{测试结果与运行记录}

\section{测试方案与验收标准}

为了让最终报告不只停留在“能运行”，本项目把测试分为环境测试、硬件测试、主线跟踪测试、鲁棒性测试和创新点测试五类。每类测试均对应固定脚本或固定操作，便于后续复现实验结果。

\begin{table}[htbp]
  \centering
  \caption{总体测试方案与验收标准}
  \label{tab:test-plan}
  \footnotesize
  \setlength{\tabcolsep}{4pt}
  \begin{tabular}{p{0.18\linewidth}p{0.36\linewidth}p{0.36\linewidth}}
    \toprule
    测试类别 & 测试内容 & 通过标准 \\
    \midrule
    环境测试 & 虚拟环境、依赖导入、单元测试、源码包导入 & \code{01\_env\_check.sh} 无错误退出，单元测试全部通过 \\
    硬件测试 & I2C、PCA9685、舵机、摄像头 & 能发现 \code{0x40}，两个舵机按安全角度动作，摄像头可读帧 \\
    主线跟踪测试 & 网页模拟模式和网页真实模式 & 浏览器能看到实时画面，目标可被框选，真实模式下云台跟随目标移动 \\
    鲁棒性测试 & 遮挡、相近颜色干扰、光照变化、目标快速移动 & 系统不乱转，短暂丢失后可重新捕获，误检不应长期持续 \\
    创新点测试 & C++ wheel、pymp、网页控制台、mock 云端上传、安全保护 & 有可运行脚本、有输出文件或截图、有可写入报告的对比数据 \\
    \bottomrule
  \end{tabular}
\end{table}

\begin{table}[htbp]
  \centering
  \caption{最终跟踪场景实测记录占位表}
  \label{tab:tracking-scenarios}
  \footnotesize
  \setlength{\tabcolsep}{4pt}
  \begin{tabular}{p{0.18\linewidth}p{0.30\linewidth}p{0.22\linewidth}p{0.20\linewidth}}
    \toprule
    场景 & 操作方式 & 当前记录 & 后续补充材料 \\
    \midrule
    静止目标 & 将目标放在画面偏左/偏右位置，观察云台是否回正 & 待补充 & 截图、平均误差 \\
    缓慢移动 & 手持目标缓慢横向移动 & 待补充 & 演示视频、误差曲线 \\
    快速移动 & 目标快速从一侧移动到另一侧 & 待补充 & 最大误差、是否丢失 \\
    短时遮挡 & 用手遮挡目标 1--2 秒后移开 & 待补充 & 重捕时间、状态切换日志 \\
    颜色干扰 & 背景放置相近颜色物体 & 待补充 & 是否误跟、目标框截图 \\
    光照变化 & 改变台灯或环境亮度 & 待补充 & HSV 调参前后截图 \\
    \bottomrule
  \end{tabular}
\end{table}

\section{基础测试结果}

当前版本已经完成的基础测试如表~\ref{tab:basic-tests} 所示。这些测试构成后续演示前的最小检查集合。

\begin{table}[htbp]
  \centering
  \caption{当前基础测试结果}
  \label{tab:basic-tests}
  \begin{tabular}{lll}
    \toprule
    测试项 & 当前结果 & 说明 \\
    \midrule
    Python 环境 & 通过 & 虚拟环境可用，依赖可导入 \\
    单元测试 & 通过 & 控制器、状态机和运行时 HSV 测试通过 \\
    I2C 扫描 & 通过 & \code{/dev/i2c-1} 可发现 \code{0x40} 和 \code{0x70} \\
    舵机测试 & 通过 & channel 0 为 tilt，channel 1 为 pan \\
    摄像头读取 & 通过 & \code{VideoCapture(0)} 可读取图像 \\
    网页模拟跟踪 & 通过 & 可显示实时画面、中心自动标定和只读 HSV 状态 \\
    网页真实跟踪 & 已跑通 & 目标检测结果可驱动真实舵机 \\
    日志分析 & 通过 & 可生成 FPS、误差、角度和目标状态曲线 \\
    \bottomrule
  \end{tabular}
\end{table}

\section{本地与真实跟踪测试}

本地模拟跟踪用于验证摄像头、视觉检测、状态机和画面叠加是否正常，但不控制真实舵机，适合在调自动标定效果和识别稳定性时使用。当前截图显示颜色目标可以被框选，并显示面积和置信度。

\begin{figure}[htbp]
  \centering
  \includegraphics[width=0.92\linewidth]{figures/local_tracking_current.jpg}
  \caption{本地模拟跟踪测试画面}
  \label{fig:local-tracking-current}
\end{figure}

真实跟踪会在识别目标的同时控制云台。当前测试中，目标能够被识别，画面叠加了目标框、中心点和中心线，终端中主程序处于跟踪运行状态。

\begin{figure}[htbp]
  \centering
  \includegraphics[width=0.92\linewidth]{figures/real_tracking_current.jpg}
  \caption{真实跟踪测试画面}
  \label{fig:real-tracking-current}
\end{figure}

\section{日志与量化指标}

程序运行时会在 \code{logs/<timestamp>} 下生成 \code{frames.csv} 和 \code{summary.json}。\code{frames.csv} 记录每帧时间戳、状态、检测耗时、控制耗时、FPS、中心误差、pan/tilt 角度、目标是否找到、面积和置信度；\code{summary.json} 记录平均帧率、平均误差、最大误差、目标丢失事件和重新捕获时间等汇总指标。

当前版本提供 \code{08\_analyze\_latest\_logs.sh} 自动分析最新日志，生成 CSV、JSON、Markdown 和曲线图。最终答辩前需要选择一段稳定运行数据，将真实曲线补入本报告。

\placeholderfigure{FPS 曲线占位}{后续运行 \code{08\_analyze\_latest\_logs.sh} 后，补入 \code{logs/<timestamp>/analysis/fps\_curve.png}。}

\placeholderfigure{中心误差曲线占位}{后续补入 \code{center\_error\_curve.png}，展示 $|e_x|$ 与 $|e_y|$ 随时间变化。}

\placeholderfigure{云台角度曲线占位}{后续补入 \code{gimbal\_angle\_curve.png}，展示 pan/tilt 角度随帧变化。}

\begin{table}[htbp]
  \centering
  \caption{最终量化指标占位表}
  \label{tab:final-metrics}
  \begin{tabular}{lll}
    \toprule
    指标 & 当前占位 & 后续来源 \\
    \midrule
    平均 FPS & 待补充 & \code{summary.json} 或 \code{tracking\_analysis.json} \\
    平均中心误差 X & 待补充 & \code{tracking\_summary\_table.csv} \\
    平均中心误差 Y & 待补充 & \code{tracking\_summary\_table.csv} \\
    最大中心误差 X/Y & 待补充 & 日志分析结果 \\
    目标丢失事件数 & 待补充 & 日志分析结果 \\
    平均重新捕获时间 & 待补充 & \code{summary.json} \\
    \bottomrule
  \end{tabular}
\end{table}

\placeholderfigure{最终跟踪日志终端输出占位}{后续补充 \code{04\_web\_console\_real.sh} 运行时的终端截图，要求能看到网页地址、运行状态和无异常退出信息。}

\placeholderfigure{最终演示视频二维码占位}{后续可将最终演示视频上传到网盘或飞书云文档，并在此处放置二维码或视频链接截图。}

\section{当前不稳定因素}

当前系统主线已经跑通，但仍有几个会影响最终演示稳定性的因素：

\begin{enumerate}
  \item Wi-Fi 热点在不同时段稳定性差异明显，可能导致 SSH 断开、网页打不开或 MJPEG 卡顿。
  \item 固定 HSV 阈值受光照影响较大，背景中出现相近颜色时可能误检。
  \item 目标快速移动或摄像头曝光不稳定时，云台会出现追踪延迟。
  \item 网页真实模式需要同时保证摄像头、I2C、舵机电源、网络和浏览器正常，任一环节异常都会影响演示。
\end{enumerate}

因此后续推荐使用 USB 转网口建立电脑与 OrangePi 的有线连接，同时由电脑共享网络给 OrangePi。这样 SSH 和网页控制台会比 Wi-Fi 热点稳定很多。

\chapter{创新点落实情况}

\section{创新点1：C/C++ wheel 视觉算子加速}

课程要求中提到可以将一个耗时功能块用 C/C++ 实现，编译成 wheel 后由 Python 调用，并与 Python 版本对比性能。当前项目选择 $5\times5$ 二值 mask 形态学开闭运算作为加速对象，原因是它位于视觉检测链路中，输入输出清晰，适合做公平 benchmark。

当前已实现四种后端：

\begin{enumerate}
  \item 纯 Python 后端，用于展示未优化循环的基准耗时。
  \item OpenCV 后端，用于主线默认运行和工程稳定性。
  \item C++/pybind11 后端，位于 \code{cpp\_accel/src/morphology\_ext.cpp}，可编译成 wheel。
  \item pymp 后端，用于对比 Python 并行化方案。
\end{enumerate}

\code{05\_cpp\_benchmark.sh} 会自动清理旧 build、构建 wheel、使用 \code{--no-deps} 安装本地 wheel，并运行 benchmark 生成 JSON、CSV 和 Markdown 结果。脚本默认使用较小图像和较少帧数，目的是先保证 OrangePi 快速跑完；正式报告数据可使用更大参数重新测试。

\begin{verbatim}
bash scripts/run/05_cpp_benchmark.sh --frames 120 --width 320 --height 240
\end{verbatim}

\placeholderfigure{C++ wheel 加速 benchmark 结果占位}{后续补充 \code{benchmark\_results/morphology\_benchmark.csv} 的表格或折线图，比较 pure Python、OpenCV、C++ wheel 和 pymp 的平均单帧耗时。}

\begin{table}[htbp]
  \centering
  \caption{C++ wheel 加速实验结果占位表}
  \label{tab:cpp-benchmark}
  \begin{tabular}{llll}
    \toprule
    后端 & 平均耗时/ms & 相对 Python 加速比 & 输出是否一致 \\
    \midrule
    pure Python & 待补充 & 1.00x & 待补充 \\
    OpenCV & 待补充 & 待补充 & 待补充 \\
    C++ wheel & 待补充 & 待补充 & 待补充 \\
    pymp & 待补充 & 待补充 & 待补充 \\
    \bottomrule
  \end{tabular}
\end{table}

\placeholderfigure{C++ wheel 编译成功截图占位}{后续补充 \code{05\_cpp\_benchmark.sh} 中 wheel 构建、安装和 benchmark 全部完成的终端截图。}

\section{创新点2：网页控制台与网络编程}

网页控制台既服务主线演示，也构成网络编程加分项。系统在 OrangePi 上启动 HTTP 服务，提供首页、MJPEG 视频流、健康检查接口和多个 JSON API。浏览器可以通过 API 读取实时状态、自动标定、开始跟踪、急停和复位。

核心接口包括：

\begin{table}[htbp]
  \centering
  \caption{网页控制台主要 API}
  \label{tab:web-api}
  \begin{tabular}{lll}
    \toprule
    接口 & 方法 & 功能 \\
    \midrule
    \code{/stream.mjpg} & GET & MJPEG 实时视频流 \\
    \code{/health} & GET & 健康检查 \\
    \code{/api/status} & GET & 获取 FPS、状态、误差、角度和 HSV 配置 \\
    \code{/api/calibrate} & POST & 根据画面中心 ROI 自动标定颜色 \\
    \code{/api/start} & POST & 开始跟踪 \\
    \code{/api/stop} & POST & 急停云台动作 \\
    \code{/api/reset} & POST & 急停并渐进复位 \\
    \bottomrule
  \end{tabular}
\end{table}

\placeholderfigure{网页控制台 API 调用结果占位}{后续补充浏览器开发者工具、curl 或脚本输出截图，展示 \code{/api/status}、\code{/api/calibrate} 或 \code{/api/start} 接口返回正常。}

\section{创新点3：pymp 并行加速对比}

Python thread 对 CPU 密集型任务通常受 GIL 限制，难以真正提升计算速度。因此项目加入 \code{pymp-pypi}，在形态学处理后端中实现了 pymp 版本，并通过 \code{06\_pymp\_benchmark.sh} 进行性能对比。

pymp 模块的意义不是替代 OpenCV 主线，而是用于回答“Python 内部如何做并行加速”这一加分问题。最终报告中可以将 pymp 与 pure Python、OpenCV、C++ wheel 放在同一张表里比较。

\placeholderfigure{pymp 并行加速结果占位}{后续补充 OrangePi Linux 环境下 pymp benchmark 结果。Windows 本地由于缺少 fork 支持，pymp 不作为有效结果。}

\section{创新点4：mock 云端服务与数据上传}

网络与云端方向目前分为两层：第一层是已经跑通的 HTTP/MJPEG 网页控制台，用于实时查看 OrangePi 处理后的画面；第二层是 mock 云端上传服务，用于把跟踪 summary 和 benchmark 结果通过 HTTP POST 发送到服务端。

\code{mock\_cloud\_server.py} 可以启动一个本地 mock 云端服务，\code{upload\_latest\_summary.py} 会读取最新一次 \code{logs/<timestamp>/summary.json}、最新 benchmark JSON，并上传到 \code{/api/tracking-runs}。一键脚本 \code{07\_mock\_cloud\_upload.sh} 会自动启动 mock 服务、打印 dry-run payload 并执行真实上传。该方案不绑定具体云平台账号，后续只需把 endpoint 替换为真实云服务器 API，即可把实验数据上传到云端。

\placeholderfigure{mock 云端上传结果占位}{后续补充 \code{07\_mock\_cloud\_upload.sh} 的终端截图，或 \code{cloud\_uploads} 目录中收到的 JSON 文件截图。}

\begin{table}[htbp]
  \centering
  \caption{mock 云端上传内容占位表}
  \label{tab:mock-cloud-fields}
  \footnotesize
  \begin{tabular}{lll}
    \toprule
    字段类别 & 示例字段 & 后续核验方式 \\
    \midrule
    跟踪摘要 & run\_id、average\_fps、lost\_count、mean\_error & 查看上传 JSON \\
    benchmark & backend、mean\_ms、speedup、consistent & 查看 benchmark JSON \\
    设备信息 & hostname、project\_path、upload\_time & 查看 mock server 记录 \\
    \bottomrule
  \end{tabular}
\end{table}

\section{创新点5：云台安全保护机制}

云台安全保护直接服务于主线跟踪。相比只追求“能动”的版本，当前系统加入了角度限位、速度限幅、死区、防抖、中心稳定保持、开始跟踪、急停、渐进复位、启动释放 PWM 和真实硬件失败显式报错等保护。这些机制可以降低舵机顶住机械结构、过冲、持续抖动和误判运行状态的风险。

\placeholderfigure{网页急停与复位功能截图占位}{后续补充点击 START、STOP 和 RESET 按钮前后的网页截图或短视频截图，用于证明安全保护不是只写在代码中。}

\section{DNN/RKNN 扩展说明}

DNN/RKNN 当前不进入主线命令流程。原因是本项目当前验收目标是高对比颜色目标跟踪，传统视觉方案更稳定、更可解释、更容易现场演示。若后续时间充足，可选择轻量检测模型，先导出 ONNX，再用 RKNN Toolkit 转换并部署到 RK3588S NPU，最终替换现有 \code{TargetDetection} 输出接口。

因此本报告只将 DNN/RKNN 作为后续扩展方向说明，不作为当前最终演示强依赖。

\placeholderfigure{DNN/RKNN 扩展占位}{若后续完成 RKNN 部署，可补充模型转换截图、板端推理截图或 FPS 对比图。}

\chapter{网络部署与最终演示方案}

\section{USB 转网口与网络共享}

项目远程调试主要依赖 SSH 和网页控制台。前期使用 Wi-Fi 热点时，白天和晚上网络质量差异明显，可能导致 MobaXterm 断开、网页无法打开、MJPEG 卡顿或延迟变大。由于轻薄本没有原生网口，后续计划购买 USB 转网口，让电脑通过有线方式直连 OrangePi，并由电脑共享网络给 OrangePi。

采用 USB 转网口后，电脑和 OrangePi 之间的数据链路不再依赖拥挤的 Wi-Fi 热点。电脑仍然可以通过 Wi-Fi 上网，同时把网络共享到有线网口；OrangePi 通过有线网口获得 IP、默认路由和 DNS，从而既能被电脑稳定 SSH 连接，也可以在需要时访问外网安装依赖或上传数据。

\placeholderfigure{USB 转网口共享网络设置截图占位}{后续补充 Windows 网络共享设置截图、OrangePi 获取到的有线 IP 截图，以及 \code{10\_network\_check.sh} 检查结果。}

\begin{table}[htbp]
  \centering
  \caption{网络部署检查项占位表}
  \label{tab:network-check}
  \footnotesize
  \begin{tabular}{lll}
    \toprule
    检查项 & 预期结果 & 后续记录 \\
    \midrule
    OrangePi 有线 IP & 能看到有线网卡 IP 地址 & 待补充 \\
    SSH 连通性 & MobaXterm 可稳定连接 & 待补充 \\
    DNS 与外网 & OrangePi 可访问外网或安装依赖 & 待补充 \\
    网页控制台端口 & 电脑浏览器可打开 \code{http://OrangePi-IP:5000} & 待补充 \\
    MJPEG 流畅度 & 视频流不明显卡顿，控制按钮响应正常 & 待补充 \\
    \bottomrule
  \end{tabular}
\end{table}

\section{最终推荐演示流程}

最终演示建议按如下顺序进行：

\begin{enumerate}
  \item 接好舵机电源、PCA9685、I2C、摄像头和 OrangePi 电源。
  \item 使用 USB 转网口连接电脑和 OrangePi，电脑开启网络共享。
  \item 在 OrangePi 终端进入 \code{~/zb/tracker\_project}。
  \item 如有需要，临时检查网络和网页端口，但不作为主线必跑步骤。
  \item 运行 \code{bash scripts/run/01\_env\_check.sh} 确认环境和单元测试通过。
  \item 运行 \code{bash scripts/run/02\_hardware\_check.sh} 确认舵机和摄像头正常。
  \item 运行 \code{bash scripts/run/03\_web\_console\_simulate.sh}，电脑浏览器确认画面、中心自动标定和只读 HSV 显示正常。
  \item 运行 \code{bash scripts/run/04\_web\_console\_real.sh}，开始真实云台跟踪演示。
  \item 演示网页自动标定、开始跟踪、急停和复位。
  \item 演示完成后运行日志分析和创新点 benchmark，补充最终量化数据。
\end{enumerate}

\section{最终需要补充的材料}

当前报告已经写入当前版本截图，后续还需要补充以下材料，让最终提交更完整：

\begin{enumerate}
  \item 最终硬件实物接线照片，最好带标注。
  \item USB 转网口共享网络成功截图。
  \item 最终远程真实跟踪演示截图或视频二维码。
  \item 一段稳定运行日志生成的 FPS、误差和角度曲线。
  \item C++ wheel / pymp benchmark 的 OrangePi 板端真实数据表格。
  \item mock 云端上传成功截图或接收到的 JSON 文件截图。
\end{enumerate}

\section{答辩展示要点}

答辩时建议按照“问题背景、硬件搭建、主线闭环、网页控制台、创新点、量化结果、后续改进”的顺序讲。这样能够避免一上来陷入代码细节，而是先让老师看到系统确实完成了从摄像头到云台动作的闭环。

\begin{table}[htbp]
  \centering
  \caption{答辩展示材料占位清单}
  \label{tab:defense-materials}
  \footnotesize
  \setlength{\tabcolsep}{4pt}
  \begin{tabular}{p{0.20\linewidth}p{0.42\linewidth}p{0.28\linewidth}}
    \toprule
    材料 & 用途 & 当前状态 \\
    \midrule
    硬件接线图 & 说明系统如何搭起来，避免硬件逻辑讲不清 & 待补最终标注图 \\
    网页控制台截图 & 展示远程监控、中心自动标定、只读 HSV 和状态面板 & 已有基础截图，待补最终版 \\
    真实跟踪视频 & 证明系统不是离线截图，而是实时闭环 & 待补最终视频或二维码 \\
    日志曲线 & 证明系统有量化评估，不只凭肉眼观察 & 待补 FPS/误差/角度曲线 \\
    加速 benchmark & 证明 C++ wheel 与 pymp 加分项有实验对比 & 待补 OrangePi 实测表 \\
    mock 云端上传截图 & 证明网络与云端方向已完成可运行原型 & 待补上传结果 \\
    \bottomrule
  \end{tabular}
\end{table}

\chapter{总结}

当前版本已经完成课程大作业的核心主线：OrangePi 能够读取摄像头图像，识别高对比颜色目标，通过 PCA9685 控制两自由度云台，并通过网页控制台展示叠加后的实时画面。项目还将调试流程整理为脚本，降低了重复测试和现场演示时的操作复杂度。

在创新点方面，项目已经落实 C++ wheel 视觉算子加速、pymp 并行对比、网络网页控制台、mock 云端上传和云台安全保护机制。这些支线不是脱离主线的“炫技”，而是服务于最终跟踪系统：C++ 和 pymp 用于解释速度优化，网页和云端用于解释网络化展示与数据同步，安全保护用于提高真实机械系统的可靠性。

后续工作的重点是把系统从“能跑”继续调到“稳跑”：使用 USB 转网口改善远程连接，调整 HSV 阈值和控制参数，补充板端 benchmark 与日志曲线，并录制一段稳定的最终演示视频。完成这些后，系统就能够形成“硬件搭建清楚、主线演示稳定、创新点可量化、答辩逻辑完整”的最终成果。

\end{document}
